\documentclass{claudeeditorial}

\renewcommand{\DocKicker}{Manual de marcha}
\renewcommand{\DocTitle}{Diario de ingeniería de proyectos}
\renewcommand{\DocSubtitle}{Plantilla diaria para decisiones, riesgos, avances, arquitectura y seguimiento operativo}
\renewcommand{\DocAuthor}{Equipo de proyecto}
\renewcommand{\DocRole}{Project engineering}
\renewcommand{\DocDate}{5 de mayo de 2026}
\renewcommand{\DocRef}{DAYBOOK-001}
\renewcommand{\DocFooter}{claudeeditorial project daybook}

\renewcommand{\ProjectName}{Sistema de operaciones}
\renewcommand{\ProjectPhase}{Ejecución}
\renewcommand{\ProjectOwner}{Equipo de ingeniería}
\renewcommand{\ProjectStatus}{En progreso}

\ClaudeRunningHeader
\setstretch{1.22}

\begin{document}
\BodyCopy

\ClaudeCover

\ClaudeChapterPage{01}{Nota diaria}{Una página de mando para que el equipo sepa qué cambió, qué se decidió y qué necesita atención.}

\ClaudeProjectHeader
  {\ProjectName}
  {Seguimiento diario de ingeniería: estado, decisiones, riesgos, evidencias y siguiente movimiento.}
  {\DocDate}
  {\ProjectStatus}

\begin{tcbraster}[raster columns=4,raster equal height=rows,raster column skip=8pt]
\KPICard{Salud}{Verde}{entregables dentro de rango}
\KPICard{Riesgos}{2}{uno técnico, uno operativo}
\KPICard{Decisiones}{1}{arquitectura de integración}
\KPICard{Bloqueos}{0}{sin dependencia externa crítica}
\end{tcbraster}

\section{Resumen ejecutivo}

\begin{claudehero}
\ClaudeLead{El proyecto avanza según el plan. La prioridad del día es cerrar el contrato de integración, dejar trazabilidad de la decisión técnica y preparar una validación reproducible para revisión de negocio.}
\end{claudehero}

\begin{ClaudeActionList}[Foco del día]
\ClaudeActionItem{Contrato de integración}{due: hoy}{Congelar nombres de campos, tipos y semántica de errores.}
\ClaudeActionItem{Validación reproducible}{due: 24 h}{Generar una consulta de humo y registrar salida esperada.}
\ClaudeActionItem{Revisión de riesgos}{due: semanal}{Confirmar que las mitigaciones siguen teniendo propietario.}
\end{ClaudeActionList}

\section{Decisión}

\ClaudeDecisionRecord[Activa]
  {ADR-004 · Contrato único de eventos}
  {Los equipos consumen eventos desde dos fuentes con campos casi equivalentes, pero con nombres y reglas de nulidad diferentes.}
  {Se define un contrato canónico y se obliga a que cada adaptador traduzca al modelo común antes de publicar.}
  {La complejidad queda en los bordes del sistema, mejora la observabilidad y se reduce el coste de soporte.}

\section{Riesgos}

\begin{ClaudeTable}{@{}lL L l@{}}
\ClaudeTableHeader Riesgo & Impacto & Mitigación & Estado \\
\toprule
\ClaudeRiskRow{R-012}{Datos incompletos en origen secundario}{Validar nulos antes de publicar eventos.}{\StatusWIP}
\ClaudeRiskRow{R-018}{Cambios de esquema sin aviso}{Añadir prueba de contrato en pipeline.}{\StatusOpen}
\bottomrule
\end{ClaudeTable}

\section{Evidencia técnica}

\begin{ClaudeCode}[title={Smoke test SQL}]{SQL}
SELECT event_type, COUNT(*) AS total
FROM ops.event_log
WHERE created_at >= CURRENT_DATE - 1
GROUP BY event_type
ORDER BY total DESC;
\end{ClaudeCode}

\begin{ClaudeConsole}
$ npm run contract:check
contract:event-log  ok
contract:payload    ok
\end{ClaudeConsole}

\section{Arquitectura}

\begin{ClaudeDiagram}
\centering
\begin{tikzpicture}[claude diagram canvas,node distance=20mm]
  \node[claude diagram node] (source) {Source};
  \node[claude diagram node,right=of source] (adapter) {Adapter};
  \node[claude diagram node,right=of adapter] (contract) {Canonical\\Contract};
  \node[claude diagram node,right=of contract] (consumer) {Consumer};
  \draw[claude uml relation] (source) -- (adapter);
  \draw[claude uml relation] (adapter) -- (contract);
  \draw[claude uml relation] (contract) -- (consumer);
\end{tikzpicture}
\end{ClaudeDiagram}

\begin{ClaudeInfo}[Cierre]
El documento diario debe ser breve. Si una decisión requiere más contexto, se abre un ADR completo y se enlaza desde esta página.
\end{ClaudeInfo}

\ClaudeSignature{\DocAuthor}{\DocRole}

\end{document}
